\documentclass[10pt]{article}
\usepackage[margin=.8in]{geometry}
\usepackage[T1]{fontenc}
\usepackage{amsmath,amssymb,booktabs,hyperref,xurl}
\hypersetup{colorlinks=true,urlcolor=blue,linkcolor=blue}
\title{Restore the Proximal CLF Objective and Diagnose Delayed Return}
\author{Collision Avoidance Research}
\date{October 3, 2026}
\begin{document}
\maketitle
\section{Corrected requirement and implementation}
The user's preference concerns lateral excursion from the given path, not
smooth actuator changes. The recently added input-continuity objective and
its configuration field are therefore removed completely. The retained
controller minimizes PCBF slack during initialization/restoration, or inherits
the shifted budget, and then minimizes the same first-successor CLF slack with
its original bounded anchor-relative input correction cost:
\[
 \min\;\frac{\rho}{\max(1,V(x))}+\epsilon R_a(\Delta U),\quad
 R_a=\frac{1}{2M}\sum_{i=0}^{M-1}\sum_{l=1}^{2}
       \frac{\Delta u_{l,i}^2}{r_l^2},\quad \epsilon=10^{-4}.
\]
The input reference is the linearization trajectory, not zero. There is no
new lateral-deviation objective, road boundary, steering magnitude/rate limit,
CLF, controller mode, backup or optimization stage. The 0.10-m affine collision
margin remains. The solver and configuration sources are byte-identical to
commit \path{b16771139356ec252ed8deb4bb341cc22ffb023a}; continuation version 70
prevents reuse of the retired objective's continuation state. Existing
configuration validation rejects the removed option explicitly.

\section{What the CLF does and what the diagnosis tests}
The unchanged function $V=e^\top P e$ uses transverse position, heading,
longitudinal/lateral velocity and yaw-rate errors. Its constraint is
\[
 V_{\mathrm{aff}}(x_{1|k})-V(x_k)
 \le-\eta e(x_k)^\top Qe(x_k)+\rho_k,\quad \eta=0.5,\quad\rho_k\ge0.
\]
It acts on the first successor. Future controls have no individual CLF
constraints; horizon feasibility and the tie objective select them. Thus a
small tie objective can affect future inputs strongly along directions that
do not change the current minimum slack. A zero-slack decrease condition is
not a minimum-excursion objective. Positive slack permits $V$ to rise, and the
local LQR quadratic is not a global nonlinear CLF.

The historical 8-m/s head-on run from smoothing commit
\path{63ccffd5a1cbf21f659c9ec7c6b167a127d44e21} was replayed from its recorded
measured states. All 161 issued inputs through 8.0 s were reproduced exactly,
with maximum input discrepancy zero. An isolated diagnostic copy captured
the actual conic problem at eight times: 1.0, 1.4, 1.45, 1.5, 2.0, 4.0,
6.0 and 8.0 s. No diagnostic candidate was executed on the plant or installed
as another production controller.

For each frozen problem, the analysis (i) restored the anchor-only objective,
(ii) removed all tie cost and minimized only CLF slack, (iii) removed target
collision/separation rows while preserving the same dynamics/CLF, and
(iv) minimized the condensed two-variable first-step CLF quadratic over only
the existing first-input correction box. The last problem relaxes all future
constraints and gives a lower bound on the original minimum slack.

\section{Root cause: future-plan selection and shifted correction bounds}
The immediate steering output is essentially unchanged under the first two
ablations: its change is below $1.3\times10^{-8}$ rad. Braking changes slightly,
but the raw CLF slack improvement stays below $3\times10^{-4}$. The relaxed
first-step minimum also agrees with the full problem within that raw slack
tolerance at all eight times. The CLF stage was being solved; neither skipping
it nor a large smoothing weight overriding it explains this regression.

Every examined frame places steering at its lower correction bound:
\[
 \delta_{0|k}=\bar\delta_{0|k}-r_{\delta,k}.
\]
This is a numerical correction limit about the planned anchor, not an actuator
or previous-command slew limit. At $t=4.0$ s, the actual stored plan is:
\begin{center}\begin{tabular}{lr}\toprule
Quantity & Steering (rad)\\\midrule
Current first-input anchor $\bar\delta_{0|k}$ & $-0.009374$\\
Allowed negative correction $-r_{\delta,k}$ & $-0.075000$\\
Issued first steering $\delta_{0|k}$ & $-0.084374$\\
Planned next steering $\delta_{1|k}$ & $-0.009398$\\
\bottomrule\end{tabular}\end{center}
Thus the current frame commands about $-4.83^\circ$, but its next planned
command is only $-0.54^\circ$. On the following frame, that small planned
command becomes the new anchor. A negative correction of at most
$4.30^\circ$ again yields approximately $-4.83^\circ$. The desired stronger
return cannot accumulate through the issued input because the update shifts
the \emph{planned next input}, not the already corrected first input.
The histories at 6 and 8 s show the same pattern.

At 1.5 s, the anchor is still $+0.077717$ rad and the current radius is
$0.037993$ rad, so even the most negative allowed correction issues
$+0.039725$ rad. The vehicle still turns away while the CLF optimizer is already
using its maximum available correction toward return. The smoothing term's
important effect is its accumulated change to these future plans and hence
to subsequent affine problems, not a large instantaneous sacrifice of CLF
priority inside an unchanged problem.

For additional evidence, an offline grid of 81 steering values in
$[-0.65,0.65]$ rad and 41 signed braking ratios in $[-0.5,0.5]$ evaluates the
true RK4 first successor. At 4 s, the best sampled input is approximately
$[-0.325,0.300]^\top$, outside the current correction box. Its required slack
is 743.988, versus 817.727 for the actual nonlinear first successor of the
issued command (826.350 in the affine prediction). Positive slack remains:
this is not proof of an available zero-slack or horizon-safe maneuver. It
shows that the existing correction box excludes stronger instantaneous
CLF improvement. The finite grid is not a global nonlinear optimum.

Removing target rows gives the same optimum at seven examined times. The
4-s target-removed diagnostic returns numerical flag $-7$ and is recorded
as unresolved, not feasible or physically impossible. The independent
first-step lower bound still matches the full problem there. No claim is
made that terminal or collision constraints never matter in other frames.

\section{Closed-loop rollback results}
All fourteen scenarios again avoid sampled collision and recover nominal
cruise. Every recorded state, successor, input, dense ODE45 state/time and
transverse error is exactly identical to the original pre-smoothing campaign.
This is a new executed campaign, not a reuse of its timing measurements.
The comparison below uses the maximum lateral projection over all 31 dense
samples per 50-ms hold, not only controller nodes.

\begin{center}\small
\begin{tabular}{rlrrrr}\toprule
 & & \multicolumn{2}{c}{Maximum offset (m)} & \multicolumn{2}{c}{Recovery (s)}\\
Speed & Scenario & Smoothing & Restored & Smoothing & Restored\\\midrule
8 & headOn & 19.019 & 5.957 & 22.40 & 12.60 \\
8 & acceleratingHeadOn & 6.402 & 5.877 & 12.70 & 12.55 \\
8 & brakingLead & 5.395 & 5.549 & 14.85 & 13.50 \\
8 & crossing & 13.176 & 5.722 & 18.40 & 12.60 \\
8 & turningCrossing & 5.603 & 5.609 & 12.05 & 12.10 \\
8 & curvedHeadOn & 17.044 & 6.086 & 27.00 & 17.90 \\
8 & curvedCrossing & 5.775 & 5.751 & 13.30 & 13.05 \\
15 & headOn & 7.181 & 7.241 & 10.20 & 10.15 \\
15 & acceleratingHeadOn & 6.869 & 6.872 & 10.30 & 10.30 \\
15 & brakingLead & 7.740 & 7.872 & 16.35 & 16.10 \\
15 & crossing & 10.816 & 10.775 & 11.85 & 11.55 \\
15 & turningCrossing & 10.369 & 10.478 & 11.55 & 11.40 \\
15 & curvedHeadOn & 11.045 & 11.145 & 65.95 & 66.00 \\
15 & curvedCrossing & 10.044 & 10.017 & 10.80 & 11.10 \\
\bottomrule\end{tabular}\end{center}
The three largest low-speed regressions are removed: head-on 19.019 to
5.957 m, curved head-on 17.044 to 6.086 m, and crossing 13.176 to 5.722 m.
The family-wide maximum is now 11.145 m in the 15-m/s curved head-on case,
at 57.873 s during its later encounter. Seven of fourteen cases have smaller
maximum excursion than the smoothing variant; other cases have small
increases. No uniform minimum-detour property is claimed.

There are 4,618 issued holds, no solver interruption, and a minimum dense
rectangle gap of 49.477 mm. The 0.10-m setting remains an affine-model
constraint rather than a continuous-plant clearance guarantee. Known target
acceleration/sideslip, exact state, no road constraints, the 8/16-step prefix
and 60-step completion are unchanged. ODE45 uses relative/absolute tolerances
$10^{-11}/10^{-12}$. Recovery requires the existing one-second post-conflict
nominal dwell; 1,500 holds is an observation cap, not a performance deadline.
No random seed, sensing error or computation delay is injected into the plant.

The measured full continuation controller call has median 14.401 ms,
95th percentile 28.553 ms and maximum 76.076 ms. The largest first call is
924.437 ms. No continuation call exceeds 100 ms; 24 exceed the 50-ms hold.
This is shared-workstation wall time, not a worst-case deadline guarantee.
There are 4,545 one-model frames, 73 fresh restarts, 4,510 one-solve frames,
33 damped frames and at most three numerical solves per frame.

All 390 focused MATLAB tests pass across 18 suites, including rejection of
the removed option, the original first-step CLF/recovery tests, and the retained
generic off-diagonal quadratic solver test. Source whitespace checks pass.
Code Analyzer retains the existing sparse-index/unused-suppression findings
and the workstation settings-file warning; no new changed-code issue appears.

The correction removes an objective that did not match the user's preference.
It does not redesign the first-step CLF or prove globally minimal lateral
excursion. The diagnosed interaction remains a consideration for future
changes to initialization, regularization or trust-region policy. No additional
lateral penalty is introduced to conceal it.
\end{document}
